\documentclass[10pt,a4paper]{amsart}
\usepackage[margin=19mm]{geometry}
\usepackage{amsmath,amssymb,mathtools}
\usepackage[scr=rsfs,cal=euler]{mathalfa}
\usepackage{xfrac}
\usepackage[unicode,hidelinks,bookmarks=false]{hyperref}
\newcommand{\ie}{i.e.\ }
\newcommand{\cf}{cf.\ }
\newcommand{\resp}{respectively\ }
\newcommand{\Subsection}{\subsection}
\providecommand{\nobreakdash}{\nobreak\hbox{-}}
\setlength{\emergencystretch}{2em}
\multlinegap0pt
\usepackage{bm}
\usepackage{bbm}
\usepackage{dsfont}
\usepackage{xspace}
\usepackage{cancel}
\usepackage{mathtools}
\usepackage{amsthm}
\makeatletter
\@ifundefined{lem}{\newtheorem{lem}{Lem}}{}
\@ifundefined{pro}{\newtheorem{pro}[lem]{Proposition}}{}
\makeatother
\DeclareMathOperator{\D}{\mathcal D}
\newcommand{\myeq}[1]{\ensuremath{\stackrel{\text{#1}}{=}}}
\title[Free $\mathbb Q$-groups are residually torsion-free nilpoten]{Free $\mathbb Q$-groups are residually torsion-free nilpotent}
\author{Andrei Jaikin-Zapirain}
\date{Source: arXiv:2212.00402v3 (CC BY 4.0)}
\begin{document}
\maketitle

\noindent\emph{Source and licence.} Free $\mathbb Q$-groups are residually torsion-free nilpotent. Version arXiv:2212.00402v3 (CC BY 4.0), distributed under the Creative Commons Attribution 4.0 International licence, as recorded in the arXiv licence metadata for this exact version and shown on the abstract page https://arxiv.org/abs/2212.00402v3. Full rights notice: licenses/arxiv-2212.00402-CC-BY-4.0.txt.\par\smallskip

\noindent\textit{Editorial scope.} Counted unit: the Pro whose source label is \texttt{carnatural} in the article (source0.tex lines 500--511, source label \texttt{carnatural}), delivered together with the article's own complete original proof of it (source0.tex lines 512--544, one \texttt{proof} environment of 33 lines). No mathematics is written, repaired or re-derived: statement and proof are the article's own text line for line. Required context retained uncounted: none. Every definition, display and label of those lines is retained unchanged. Omitted from this package and not redistributed: 1--499, 545--1335. Nothing inside a retained excerpt is omitted or altered: every retained body is delivered exactly as the article's lines are written, so no omission receipt of any kind accompanies this package. Citations: the retained excerpts carry no citation command at all, so no bibliography entry of the article is transcribed and no reference is redistributed. Reference adaptation: none was needed and none was made; every reference inside the delivered unit resolves inside it, so no unresolved reference or citation is produced. Printed slips kept literal: no printed word, symbol, hypothesis or number of the counted statement or of its proof was altered. Layout and class: the article's own class is \texttt{amsart} with packages including amsmath, amssymb, amsthm, hyperref, amsfonts, mathtools, thmtools; the wrapper is the lane's pinned \texttt{amsart} editorial wrapper at 10pt on A4, which restates the theorem-like environments the retained text needs and adds the article's own macro definitions that the retained text needs (\texttt{\textbackslash D}, \texttt{\textbackslash myeq}), with extra wrapper packages (bm, bbm, dsfont, xspace, cancel, mathtools). The delivered numbering is the unit's own. Assets: no retained span contains a figure environment, a graphics-inclusion command, a PDF-page inclusion, a table or a TikZ picture; no image file is copied into this package, the package declares no asset dependency and no image file is redistributed with it. Licence: version arXiv:2212.00402v3 of this article is distributed under the Creative Commons Attribution 4.0 International licence, as recorded in the arXiv licence metadata for this exact version and shown on the abstract page https://arxiv.org/abs/2212.00402v3. A full rights notice is delivered at licenses/arxiv-2212.00402-CC-BY-4.0.txt. Characters: every retained excerpt is pure ASCII, so the delivered engine, XeLaTeX with its default Latin Modern text font, reports no missing character.
\begin{pro}\label{carnatural} Let $R=S*G$ be a crossed product with $G$ finite and let $R\hookrightarrow \D$ be a division  $R$-ring of fractions. Denote by $\D_e$ the division closure of $S$ in $\D$. Then the following are equivalent.

\begin{enumerate}
\item[(a)] $\widetilde{\dim_{\D_e}}=\dim_{\D}$ as Sylvester functions on $R$.
\item [(b)] $\dim_{\D_e} \D= |G|$.
\item[(c)]$\D$ is isomorphic to a crossed product $\D_e*G$.
\item[(d)] $\D$ is isomorphic to $\D_e\otimes_ S R$ as $(\D_e , R)$-bimodule.
\item [(e)]  $\D$ is isomorphic to $R\otimes_ S \D_e$ as $(  R,\D_e)$-bimodule.
\end{enumerate}


\end{pro}

\begin{proof} For any $h,g\in G$, define $\alpha(g,h)=u_gu_h(u_{gh})^{-1}\in R$. Since the conjugation by $u_g$ fixes $S$, it also fixes $\D_e$. Therefore, we can define a ring structure on $T=\oplus_{g\in G}\D_ev_g$ defining the multiplication on homogenous elements by means of
$$ (d_1v_g)\cdot (d_2v_h)=(d_1u_gd_2(u_g)^{-1}\alpha(g,h))v_{gh}, \ d_1,d_2\in \D_e, \ g,h\in G.$$
It is clear that $T=\D_e*G$ is a crossed product, it contains  $R$ as a subring, and $T$ is isomorphic to $\D_e\otimes_ S R$ as $(\D_e , R)$-bimodule.

There exists a natural map $\gamma: T\to \D$, sending $\sum_{g\in G} d_gv_g$ ($d_g\in D_e$) to $\sum_{g\in G} d_gu_g$. Observe that $\gamma(T)$ is a domain and of finite dimension over $\D_e$. Thus, $\gamma(T)$ is a division subring of $\D$. Since $T$ contains $R$, $\D=\gamma(T)$.  This implies that (c) and (d) are equivalent. 

Since $\dim_{D_e} T=|G|$,  (b) implies that $\gamma$ is an isomorphism, and so, (b) implies (c).

Now, let us assume (d). Let $M$ be a finitely presented  $R$-module. We have the following.
\begin{multline*}
\widetilde{\dim_{\D_e}}\ M= \frac{  \dim_{\D_e}(\D_e\otimes _S M)}{|G|}=
\frac{\dim_{\D_e}(\D_e\otimes_S(R\otimes _R M))}{|G|}=\\
\frac{\dim_{\D_e}((\D_e\otimes_SR)\otimes _R M))}{|G|}\myeq{(d)}
\frac{\dim_{\D_e}(\D\otimes_R M)}{|G|}=\\ \dim_{\D} (\D\otimes_R M)=\dim_{\D} M.\end{multline*}
This proves (a).

Now, we assume that (a) holds. Since $\D_e$ es an epic $S$-ring $\D_e\otimes_S \D_e$ is isomorphic to $\D_e$ as $\D_e$-bimodule and by the same reason, $\D \otimes_R \D $ is isomorphic to $\D $ as $\D $-bimodule. Consider $M=\D$ as an $R$-module and $N=\D_e$ as a $S$-module. Then 
\begin{multline*}1=\dim_{\D}(\D\otimes_R M)=\dim_{\D} M\myeq{(a)}\widetilde{\dim_{\D_e}} M=\\
\frac{\dim_{\D_e}(\D_e\otimes_S M)}{|G|}=\frac{\dim_{\D_e}(\D_e\otimes_S(N^{\dim_{\D_e} \D}))}{|G|}=\frac{\dim_{\D_e} \D}{|G|}
\end{multline*}
This implies (b).


Let $R^{op}$ denotes the opposite ring, that is the ring with the same elements and addition operation, but with the multiplication performed in the reverse order. Then $R^{op}\cong S^{op}*G$ and the condition (c) is equivalent to
\begin{enumerate}
 \item[(c')] $\D^{op}$ is isomorphic to a crossed product $(\D_e)^{op}*G$.
 \end{enumerate}
Our previous proof gives that (c') is equivalent to  
\begin{enumerate}
 \item[(d')] $\D^{op}$ is isomorphic to $(\D_e)^{op}\otimes_ {S^{op}}R^{op}$ as $((\D_e)^{op} , R^{op})$-bimodule.
  \end{enumerate}
Now, observe  that (d') is equivalent to (e).
\end{proof}

\end{document}
